\documentclass[11pt,a4paper]{article}

\usepackage[T1]{fontenc}
\usepackage[utf8]{inputenc}
\usepackage{lmodern}
\usepackage[a4paper,margin=25mm]{geometry}
\usepackage{microtype}
\usepackage{booktabs}
\usepackage{tabularx}
\usepackage{longtable}
\usepackage{array}
\usepackage{enumitem}
\usepackage{xcolor}
\usepackage{listings}
\usepackage{hyperref}
\usepackage{fancyhdr}

\hypersetup{
  colorlinks=true,
  linkcolor=blue!50!black,
  urlcolor=blue!60!black,
  citecolor=blue!50!black,
  pdftitle={Assignment 2: Connect Four AI Agent},
  pdfauthor={Connect Four Assignment Team}
}

\pagestyle{fancy}
\fancyhf{}
\lhead{Artificial Intelligence -- Assignment 2}
\rhead{Connect Four AI Agent}
\cfoot{\thepage}
\setlength{\headheight}{14pt}
\setlist[itemize]{leftmargin=*,itemsep=2pt,topsep=4pt}
\setlist[enumerate]{leftmargin=*,itemsep=2pt,topsep=4pt}
\renewcommand{\arraystretch}{1.15}

\definecolor{codebg}{RGB}{246,247,249}
\definecolor{codeframe}{RGB}{210,214,220}
\lstdefinestyle{projectcode}{
  basicstyle=\ttfamily\small,
  backgroundcolor=\color{codebg},
  frame=single,
  rulecolor=\color{codeframe},
  columns=fullflexible,
  keepspaces=true,
  breaklines=true,
  showstringspaces=false,
  tabsize=4
}
\lstset{style=projectcode}

\title{\textbf{Assignment 2: Connect Four AI Agent}\\
\large Minimax Search with Alpha--Beta Pruning}
\author{Alexandria University\\
Faculty of Engineering\\
Computer and Systems Department\\[8pt]
Student 1:Eyad Amr Asaad Abdo Saeed 23010312 \\
Student 2:Hany Ayman Mahmoud 23010940 \\
}
\date{}
\begin{document}
\maketitle
\thispagestyle{empty}


\tableofcontents
\newpage

\section{Assignment Requirements and Completion}

The assignment requires a complete human-versus-computer game using Minimax
with alpha--beta pruning, a search cutoff at a user-selected $K$, a heuristic
that evaluates cutoff boards, and a readable tree printed after computer
searches. The game must continue until the board is full. A GUI and GUI tree
display are bonus features. The report is also a required deliverable and must
include sample runs and trees, performance analysis at different $K$ values,
algorithms, data structures, assumptions, and other relevant details.

\begin{longtable}{@{}>{\raggedright\arraybackslash}p{0.25\textwidth}@{\hspace{5pt}}>{\raggedright\arraybackslash}p{0.12\textwidth}@{\hspace{5pt}}>{\raggedright\arraybackslash}p{0.57\textwidth}@{}}
\toprule
\textbf{Requirement} & \textbf{Status} & \textbf{Project implementation}\\
\midrule
Human plays against the computer & Complete & The console game and GUI provide a human-versus-AI match.\\
Minimax with alpha--beta pruning & Complete & \texttt{algorithms/alpha\_beta.py} searches from the computer's perspective, with MAX representing the computer and MIN representing the human.\\
Continue until the board is full & Complete & Four-in-a-row does not stop the match. Terminal state means the board is full. The final score counts all length-four windows for each owner.\\
Board at least 7 columns by 6 rows & Complete & The default board is 7 columns by 6 rows; setup validation enforces those minimum dimensions.\\
User-selected depth $K$ & Complete & The terminal runner prompts for $K$ before play. The GUI setup screen also collects search settings.\\
Heuristic at cutoff depth & Complete & Three evaluators are available: scoring opportunities, gravity-aware windows, and weighted windows.\\
Readable search tree printed to console & Complete & The console game prints the recorded alpha--beta tree on every computer turn, including leaf heuristic values, explored branches, bounds, and pruned columns.\\
GUI bonus & Complete & A PySide6 multi-screen game is included.\\
GUI tree bonus & Complete & The GUI includes a tree canvas, search playback, and alpha--beta trace and computation tables.\\
Performance analysis at multiple $K$ & Complete & The reproducible analysis script and CSV contain measurements for $K=1$ through $K=6$.\\
Written report & This document & Complete after student names and group information are filled in and the report is reviewed.\\
\bottomrule
\end{longtable}

\section{Game Rules and Assumptions}

\subsection{Board and turns}
The default game uses 6 rows and 7 columns. A move selects a non-full column;
the disc falls to its lowest empty cell. The board is rendered with \texttt{X}
for the human, \texttt{O} for the computer, and \texttt{.} for an empty cell.
Columns are displayed to users starting at 1. Internally, columns are indexed
from 0.

\subsection{Scoring and ending}
The game does not end when a player first makes four in a row. Play continues
until all 42 cells are occupied. Each contiguous horizontal, vertical, or
diagonal window of exactly four same-owner discs contributes one point. Windows
overlap independently: a run of five discs contains two connected-four windows.
At the end, the computer utility is
\[
  U(s)=F_{\mathrm{computer}}(s)-F_{\mathrm{human}}(s),
\]
where $F$ counts all completed four-cell windows. A positive utility means a
computer win, a negative utility means a human win, and zero is a draw.

\subsection{Search assumptions}
\begin{itemize}
  \item The game is deterministic, two-player, zero-sum, and has perfect information.
  \item The computer maximizes the utility and the human is modeled as minimizing it.
  \item The same fixed perspective (computer minus human) is used at every node.
  \item Search stops at depth $K$ or earlier if the board becomes full. A full-board
        utility is exact and takes precedence over heuristic evaluation.
  \item The current heuristic implementations evaluate connect-four windows and
        therefore require a connect length of four.
\end{itemize}

\section{Software Architecture}

The project separates game rules, search, evaluation, recording, setup, and user
interfaces. The main dependency flow is:

\begin{lstlisting}
main.py
  +-- ui/                 PySide6 application and screens (default)
  +-- game/console_game.py Terminal game (--console)

core/                     Board, discs, rules, game state
algorithms/               Alpha-beta search and move ordering
heuristics/               Three cutoff evaluators
tree_essentials/           Search tree, recording, statistics, printing
setup/ and options/        Settings validation and game construction
experiments/               Reproducible comparisons and analysis
tests/                     Unit tests
\end{lstlisting}

\subsection{Important packages and classes}
\begin{description}[style=nextline,leftmargin=28mm]
  \item[\texttt{core/board.py}] \texttt{Board} stores an immutable grid and
  implements legal columns, landing rows, and disc drops.
  \item[\texttt{core/disc.py}] \texttt{Disc} stores owner and display color.
  Shared instances are reused for equal owner/color pairs.
  \item[\texttt{core/rules.py}] \texttt{Rules} builds and stores all possible
  connect-four windows and counts completed windows.
  \item[\texttt{core/game.py}] \texttt{GameState} pairs a board with the side
  to move. \texttt{Game} applies moves, alternates turns, checks full-board
  termination, and calculates utility.
  \item[\texttt{algorithms/alpha\_beta.py}] \texttt{AlphaBetaSearcher} runs
  recursive depth-limited Minimax with alpha--beta pruning.
  \item[\texttt{algorithms/ordering.py}] Composable ranking rules prioritize
  tactically useful columns and improve pruning efficiency.
  \item[\texttt{heuristics/}] The evaluator implementations and weight data
  structures are described in Section~\ref{sec:heuristics}.
  \item[\texttt{tree\_essentials/node.py}] \texttt{TreeNode} stores a search
  position's role, depth, value, alpha/beta window, child nodes, best column,
  and pruned columns.
  \item[\texttt{tree\_essentials/search\_recorder.py}] A recorder interface
  separates the search algorithm from output. \texttt{TreeRecorder} stores a
  tree; \texttt{NullRecorder} avoids tree-storage overhead during experiments.
  \item[\texttt{tree\_essentials/search\_stats.py}] \texttt{SearchStats}
  tracks visited and expanded nodes, heuristic evaluations, terminal leaves,
  cutoffs, pruned branches, maximum depth, and elapsed time.
  \item[\texttt{tree\_essentials/tree\_printer.py}] \texttt{TreePrinter}
  formats search results and trees for console output or export.
  \item[\texttt{setup/}] Settings validation, player construction, session
  creation, and turn history.
  \item[\texttt{ui/}] PySide6 menu, setup, battle, results, archive, tree
  visualization, and alpha--beta playback/trace views.
\end{description}

\section{Minimax with Alpha--Beta Pruning}

At each node the search generates legal columns, orders them, applies each
candidate move, and searches the resulting state. At a MAX node it retains the
largest child value and updates $\alpha$. At a MIN node it retains the smallest
child value and updates $\beta$. A branch is skipped when $\alpha\geq\beta$;
the skipped columns are saved in the recorded tree.

\begin{lstlisting}
ALPHA_BETA(state, depth, alpha, beta, role):
    if board is full:
        return exact final utility
    if depth == K:
        return heuristic(board)
    if role == MAX:
        value = -infinity
        for move in ordered_legal_moves(state):
            value = max(value, ALPHA_BETA(child(move), depth+1,
                                         alpha, beta, MIN))
            alpha = max(alpha, value)
            if alpha >= beta: record skipped moves; break
    else:
        value = +infinity
        for move in ordered_legal_moves(state):
            value = min(value, ALPHA_BETA(child(move), depth+1,
                                         alpha, beta, MAX))
            beta = min(beta, value)
            if beta <= alpha: record skipped moves; break
    return value
\end{lstlisting}

The search records whether node values are exact or bounds when pruning means
that a node was not fully explored. Alpha and beta are stored on internal tree
nodes. At each computer turn the AI performs a new search and plays the best
root column. The tree recorder is enabled for the game. The analysis program
disables recording to avoid including its time and memory cost in performance
measurements.

\subsection{Move ordering}
Move ordering affects which branch is searched first and therefore how much
alpha--beta pruning occurs; it does not change the Minimax result. The default
order is \texttt{four\_gain\_center}: first prioritize moves that create more
completed fours, then prefer columns closer to the center. Other implemented
choices include center-first, threat-first with center tie-breaking, a combined
tactical order, and left-to-right order. The tactical ordering estimates
immediate four gains and opponent reply opportunities.

\section{Heuristic Evaluators}\label{sec:heuristics}

All evaluators return a score from the computer's perspective. For raw feature
score $r$, the shared normalization is
\[
  N(r)=\frac{r}{1000+|r|},
\]
which preserves score ordering and maps finite raw values into $(-1,1)$.

\subsection{Weighted windows (Hard)}
For each four-cell window, the evaluator ignores windows containing both
players' discs. It counts completed fours, open threes (three discs and one
empty cell), open twos (two discs and two empty cells), and center-column
discs. For the current default weights,
\[
 r_{\mathrm{windows}}=
 100\Delta F + 10\Delta O_3 + 5\Delta O_2 + 10\Delta C,
\]
where each $\Delta$ is the computer's count minus the human's count. On a
seven-column board the center feature counts the central column. This is the
current Hard evaluator and is also the default evaluator in the GUI.

\subsection{Gravity-aware windows (Intermediate)}
This evaluator uses the same types of windows but discounts open threes and
twos whose empty target cells have empty cells beneath them. Such cells must be
filled before a disc can land at a target. If $s$ distinct support cells are
needed, that feature contributes $0.5^s$ instead of a full count. Completed
fours are not discounted. It currently uses the baseline weights
$(100,10,3,2)$ for fours, open threes, open twos, and center discs.

\subsection{Scoring opportunities (Easy)}
For each player and each currently playable column, the evaluator finds the
landing cell and checks only the four-cell windows containing that cell. If
the other three cells in a window already belong to that player, playing the
column would complete a four. It counts (1) columns with at least one such
gain, (2) the largest gain in any one column, and (3) the total gain across
columns. With the default parameters,
\[
 r_{\mathrm{opportunities}}=
 100\Delta F + 20\Delta M + 15\Delta G_{\max}
 + 5\Delta G_{\mathrm{total}} + 2\Delta C.
\]
This model focuses on immediate scoring options rather than valuing all open
patterns.

\subsection{Heuristic difficulty models}
The console difficulty selector maps Easy to scoring opportunities,
Intermediate to gravity-aware windows, and Hard to the weighted-window model.
The GUI presents the evaluators through its three AI profiles. Difficulty
labels describe the intended model order, not a guarantee of win rate against
every person or at every value of $K$.

\section{Weight Tuning and Heuristic Comparison}

The weighted-window evaluator's four weight remains 100 and the open-three
weight remains 10. Earlier hill-climbing and genetic-algorithm experiments
compared the candidate weights against the original baseline evaluator.
The genetic candidate $(100,10,5,6)$ performed better overall than the
hill-climbing candidate $(100,10,2,4)$ on the recorded validation set. The
center-control weight reached the top of the original search range, so a
follow-up sweep tested center weights from 0 through 20 while holding open
threes at 10 and open twos at 5.

The focused validation compared center weights 6, 10, 14, 16, and 18 against
the original baseline, with 64 games per candidate at each search depth. The
table gives each candidate's average final score margin (candidate score
minus baseline score), with positive values favoring the candidate.

\begin{table}[ht]
\centering
\caption{Center-weight validation. Open-three weight = 10; open-two weight = 5.}
\begin{tabular}{crrrr}
\toprule
Center weight & $K=2$ & $K=3$ & $K=4$ & Mean over $K$\\
\midrule
6  & +0.344 & +0.750 & +0.469 & +0.521\\
10 & +0.500 & +0.563 & +1.156 & +0.740\\
14 & +0.594 & +0.438 & +1.000 & +0.677\\
16 & +0.406 & +0.000 & +0.750 & +0.385\\
18 & +0.406 & +0.000 & +0.719 & +0.375\\
\bottomrule
\end{tabular}
\end{table}

The selected center weight is 10 because it had the highest mean margin over
the three tested depths. The depth-specific winner differs, which shows that
weight quality depends partly on search depth and test positions. These are
empirical self-play comparisons against the original evaluator, not direct
measurements of difficulty against human players. The complete CSV records are
\texttt{experiments/heuristic\_tuning.csv},
\texttt{experiments/center\_weight\_tuning.csv}, and
\texttt{experiments/center\_weight\_validation.csv}.

\section{Search-Cost Experiment}

The analysis script uses the same reproducible positions at every tested
depth. The final saved run uses 20 positions generated from legal random
openings of five plies, seed 20261125, the current window evaluator with
weights $(100,10,5,10)$, and the four-gain-center ordering. Tree recording is
off during timing. Values below
are means over the 20 positions; timings are machine- and runtime-dependent.

\begin{table}[ht]
\centering
\caption{Alpha--beta work and timing on the saved 20-position run.}
\begin{tabular}{rrrrrr}
\toprule
$K$ & Visited & Expanded & Pruned branches & Mean time (ms) & Median time (ms)\\
\midrule
1 & 8.0 & 1.0 & 0.0 & 1.096 & 1.076\\
2 & 21.6 & 8.0 & 35.4 & 1.578 & 1.493\\
3 & 85.4 & 23.6 & 80.8 & 5.375 & 4.301\\
4 & 218.8 & 87.0 & 391.6 & 9.348 & 9.103\\
5 & 760.0 & 232.6 & 869.0 & 32.618 & 30.515\\
6 & 1878.1 & 755.6 & 3404.4 & 73.005 & 69.226\\
\bottomrule
\end{tabular}
\end{table}

The measurements show the expected increase in work as $K$ grows, despite
pruning. The large pruned-branch count is the number of legal child moves
skipped at cutoffs, summed across the search; it is not the number of nodes
that were visited. Search times are averages over small searches from
mid-opening positions and should not be interpreted as a worst-case latency
for every turn.

The final results are saved in \texttt{k\_analysis\_final.csv} in the
\texttt{experiments} directory. The earlier K-analysis file predates the
final center weight and is retained as a historical run. The experiment can
be repeated with:
\begin{lstlisting}
python -m experiments.k_analysis \
  --ks 1 2 3 4 5 6 --positions 20 --opening-plies 5 \
  --csv experiments/k_analysis_final.csv
\end{lstlisting}

\section{Representative Sample Run and Search Tree}

The following trace is generated by the actual alpha--beta implementation from
a fixed legal opening. It uses the tuned weighted-window evaluator, default
four-gain-center ordering, and $K=2$. The opening is: human column 1, computer
column 7, human column 2, computer column 7, human column 3. The computer then
searches from the resulting position. Column labels in the tree are zero-based,
as in the implementation; the chosen internal column 3 is the user's column 4.

\begin{lstlisting}
Search result
  algorithm : alpha-beta
  K         : 2
  best col  : 3
  value     : +0.0050
  visited   : 21   expanded: 8   evaluations: 13
  cutoffs   : 6   pruned branches: 36
  time      : approximately 0.0010 s

ROOT [MAX] value=+0.0050 (a=-inf, b=+inf)
|-- * col 3 [MIN] value=+0.0050 (a=-inf, b=+inf)
|   |-- * col 3 LEAF eval=+0.0050
|   |-- col 2 LEAF eval=+0.0099
|   |-- col 4 LEAF eval=+0.0148
|   |-- col 1 LEAF eval=+0.0099
|   |-- col 5 LEAF eval=+0.0148
|   |-- col 0 LEAF eval=+0.0148
|   \-- col 6 LEAF eval=+0.0148
|-- col 2 [MIN] value<=-0.1071 (a=+0.0050, b=+inf)
|   |-- * col 3 LEAF eval=-0.1071
|   \-- x pruned columns: 2, 4, 1, 5, 0, 6
|-- col 4 [MIN] value<=-0.0950 (a=+0.0050, b=+inf)
|   |-- * col 3 LEAF eval=-0.0950
|   \-- x pruned columns: 2, 4, 1, 5, 0, 6
|-- col 1 [MIN] value<=-0.1071 (a=+0.0050, b=+inf)
|   |-- * col 3 LEAF eval=-0.1071
|   \-- x pruned columns: 2, 4, 1, 5, 0, 6
|-- col 5 [MIN] value<=-0.1031 (a=+0.0050, b=+inf)
|   |-- * col 3 LEAF eval=-0.1031
|   \-- x pruned columns: 2, 4, 1, 5, 0, 6
|-- col 0 [MIN] value<=-0.1071 (a=+0.0050, b=+inf)
|   |-- * col 3 LEAF eval=-0.1071
|   \-- x pruned columns: 2, 4, 1, 5, 0, 6
\-- col 6 [MIN] value<=-0.0991 (a=+0.0050, b=+inf)
    |-- * col 3 LEAF eval=-0.0991
    \-- x pruned columns: 2, 4, 1, 5, 0, 6
\end{lstlisting}

The root is MAX, so the computer chooses the child with the greatest backed-up
value. Each MIN child represents a possible human reply. For the first
computer candidate, all seven replies were evaluated. For later candidates,
the first human reply was already sufficient to prove an upper bound no better
than the root's current alpha; the other six replies were pruned. Thus the
reported tree contains all visited nodes and explicitly lists skipped moves.
The standalone console printer normally formats values to two decimals. The
sample above shows four decimal places so that small early-game heuristic
differences are visible; the GUI's detailed trace also displays four decimals
and raw feature scores.

\subsection{Interpreting a leaf value}
For example, the leaf value $-0.1071$ is an estimate from the computer's
perspective. Its negative sign means the heuristic considers that cutoff board
favorable to the human. It is not a guaranteed final score: the board is not
full, and the search stopped at depth $K=2$.

\section{Graphical Interface and Additional Features}

The GUI is launched by default from \texttt{main.py}. It includes a themed
main menu, story and map navigation, opponent/difficulty selection, match
configuration, a battle screen, a result summary, and an archive for
experiments and analysis. The search runs in a background Qt worker so that
the interface can remain responsive while the AI searches.

The tree view displays MAX/MIN nodes, backed-up values, alpha/beta bounds,
chosen moves, and pruned branches. It supports expansion, selection, zoom and
pan, and step-by-step playback. The trace tools derive a visiting-order
alpha--beta table and a bottom-up Minimax computation table from the recorded
tree. Evaluator inspection exposes component features, normalized values,
and raw weighted scores.

The GUI trace is enabled by default and has a visible control. For the
assignment's guaranteed full console tree on every AI move, use the terminal
mode described in Section~\ref{sec:run}. GUI console reports are intentionally
bounded in line count on large searches; the interactive GUI tree is the place
to inspect the recorded tree in detail.

\section{How to Open, Install, and Run}\label{sec:run}

\subsection{Open the project}
Open the project root in an IDE such as VS Code or PyCharm, or open a terminal
there. The project root is the directory containing \texttt{main.py},
\texttt{README.md}, \texttt{requirements.txt}, and the package directories.
Commands below assume the current directory is the project root.

\subsection{Linux or macOS}
Python 3.10 or newer is recommended. PySide6 is installed from the project's
requirements file.
\begin{lstlisting}
cd /path/to/Connect-4
python3 -m venv .venv
source .venv/bin/activate
python -m pip install --upgrade pip
python -m pip install -r requirements.txt
\end{lstlisting}

Launch the GUI or terminal game:
\begin{lstlisting}
python main.py                 # GUI (default)
python main.py --console       # Terminal human-versus-computer game
\end{lstlisting}

The console game asks for a positive integer $K$ and an AI difficulty. Enter a
column from 1 through 7 when prompted. The search tree is printed after each
computer search. The game ends only after the board is full.

\subsection{Windows}
From PowerShell in the project root:
\begin{lstlisting}
py -m venv .venv
.venv\\Scripts\\Activate.ps1
python -m pip install --upgrade pip
python -m pip install -r requirements.txt
python main.py
python main.py --console
\end{lstlisting}

If PowerShell blocks activation, the virtual environment's interpreter can be
called directly as \texttt{.venv\\Scripts\\python.exe}.

\subsection{Tests and experiments}
Run the full test suite from the project root with the virtual environment
active:
\begin{lstlisting}
python -m unittest discover tests
\end{lstlisting}
The verified project run reports 142 tests passing. To repeat the search cost
analysis, compare heuristics and orderings, or run weight tuning:
\begin{lstlisting}
python -m experiments.k_analysis --ks 1 2 3 4 5 6 \
  --positions 20 --opening-plies 5

python -m experiments.compare_heuristics \
  --mode all --ks 2 3 4 --rounds 10 \
  --csv experiments/results.csv

python -m experiments.tune_heuristic \
  --budget 12 --rounds 1 --validation-rounds 5
\end{lstlisting}
Experiment output is written to CSV files in \texttt{experiments/}.

To compile this report from the project root, run LaTeX twice so the table of
contents and page references are updated:
\begin{lstlisting}
pdflatex -interaction=nonstopmode -output-directory=report \
  report/assignment2_report.tex
pdflatex -interaction=nonstopmode -output-directory=report \
  report/assignment2_report.tex
\end{lstlisting}

\section{Conclusion}

The project implements the assignment's complete game loop and required
alpha--beta search, heuristic cutoff, configurable depth, and console search
tree. It also includes three heuristic models, move-ordering alternatives,
search and tuning experiments, a graphical game, and a GUI search-tree/trace
viewer. The saved measurements demonstrate the cost increase as search depth
grows and provide reproducible evidence for the current evaluator weights.
Before submission, fill in student and group details, verify the final
repository state, and make sure this report's experimental files correspond to
the submitted code.

\appendix
\section{Project File Map}
\begin{lstlisting}
main.py                       Application entry point
core/                         Board, discs, rules, game state
algorithms/                   Alpha-beta search and move ordering
heuristics/                   Window, gravity, opportunity evaluators
tree_essentials/              Tree nodes, recorder, printer, statistics
game/                         Human/AI players and console game
options/                      Board/search/color/first-player options
setup/                        Settings, validation, player factory, session
ui/                           PySide6 screens, tree canvas, trace playback
assets/                       GUI artwork and extracted image assets
experiments/                  Comparisons, tuning, K-analysis scripts and CSVs
tests/                        Unit tests and fixed-tree fixtures
requirements.txt              Python runtime dependency list
README.md                     Project overview and quick-start guide
report/assignment2_report.tex This report source
report/assignment2_report.pdf Compiled reading copy of this report
experiments/k_analysis_final.csv Current-weight K-cost measurements
\end{lstlisting}
\end{enumerate}

\end{document}
